\documentclass[10pt,a4paper]{amsart}
\usepackage[margin=19mm]{geometry}
\usepackage{amsmath,amssymb,mathtools}
\usepackage[scr=rsfs,cal=euler]{mathalfa}
\usepackage{xfrac}
\usepackage[unicode,hidelinks,bookmarks=false]{hyperref}
\newcommand{\ie}{i.e.\ }
\newcommand{\cf}{cf.\ }
\newcommand{\resp}{respectively\ }
\newcommand{\Subsection}{\subsection}
\providecommand{\nobreakdash}{\nobreak\hbox{-}}
\setlength{\emergencystretch}{2em}
\multlinegap0pt
\newtheorem{theorem}{Theorem}
\newtheorem{lemma}[theorem]{Lemma}
\newcommand{\NN}{\mathbb Z_{\geq0}}
\hypersetup{pdfauthor={Mubin Shaikh},
pdftitle={Rationality and Quasipolynomiality of Restricted Rectangle Partitions}}
\begin{document}
\title[Rationality and Quasipolynomiality of Restricted Rectangle P]{Rationality and Quasipolynomiality of Restricted Rectangle Partitions}
\author{Mubin Shaikh; Independent Researcher; ORCID: https://orcid.org/0009-0005-1290-7114 https://orcid.org/0009-0005-1290-7114}
\date{}
\maketitle
\noindent\emph{Source and licence.} Rationality and Quasipolynomiality of Restricted Rectangle Partitions. Version arXiv:2609.13305v1 (CC BY 4.0). This material is distributed under the Creative Commons Attribution 4.0 International licence, http://creativecommons.org/licenses/by/4.0/, as recorded in the arXiv OAI licence metadata for this exact version and shown on the abstract page https://arxiv.org/abs/2609.13305v1. Full rights notice: licenses/arxiv-2609.13305-CC-BY-4.0.txt.\par\smallskip
\noindent\textit{Editorial scope.} Counted unit: the original \texttt{lemma} environment \texttt{lem:series} at source lines 92--105 together with its complete proof at lines 106--118; the delivered document keeps source lines 61--119 so that every referenced label and every citation resolves inside it. Native editable source is the paper's own concatenated TeX; the AMS wrapper is editorial formatting only. No publisher class, upstream script or unrelated artwork is redistributed.
\section{Finite generation of upper sets}
Equip $\NN^d$ with the coordinatewise order. The next elementary form
of Dickson's lemma is included to make the argument self-contained.
\begin{lemma}\label{lem:dickson}
Every subset of $\NN^d$ has finitely many coordinatewise minimal
elements. If $U\subseteq\NN^d$ is upward closed, then
\[
 U=\bigcup_{a\in M}(a+\NN^d),
\]
where $M$ is its finite set of minimal elements.
\end{lemma}
\begin{proof}
Every infinite sequence of nonnegative integers has an infinite
nondecreasing subsequence. Indeed, either one value occurs infinitely
often, or every bounded set contains only finitely many terms and a
strictly increasing subsequence can be selected successively.
Apply this observation successively to each coordinate of an infinite
sequence in $\NN^d$. The resulting infinite subsequence is nondecreasing
in every coordinate. Thus an infinite set of distinct mutually
incomparable elements is impossible. Distinct minimal elements are
incomparable, proving finiteness.

For $u\in U$, the nonempty finite set
$\{v\in U:v\leq u\}$ contains a minimal element $a$. Any element of
$U$ strictly below $a$ would also be below $u$, so $a$ is minimal in
$U$ itself. Hence $u\in a+\NN^d$. The reverse inclusion follows
from upward closure. The empty set is covered by taking $M=\varnothing$.
\end{proof}

For positive integer weights $w=(w_1,\ldots,w_d)$, put
$w\cdot a=\sum_jw_ja_j$.

\begin{lemma}\label{lem:series}
If $U\subseteq\NN^d$ is upward closed, then
\[
 \sum_{u\in U}x^{w\cdot u}
 =\frac{A_U(x)}{\prod_{j=1}^d(1-x^{w_j})}
 \quad\text{for some }A_U(x)\in\mathbb Z[x].
\]
More explicitly, for the finite minimal set $M$,
\begin{equation}\label{eq:ie}
 A_U(x)=\sum_{\varnothing\ne S\subseteq M}
 (-1)^{|S|+1}x^{w\cdot\bigvee S},
\end{equation}
where $\bigvee S$ is the coordinatewise maximum of the vectors in $S$.
\end{lemma}

\begin{proof}
For each nonempty $S\subseteq M$,
\[
 \bigcap_{a\in S}(a+\NN^d)=(\bigvee S)+\NN^d,
 \qquad
 \sum_{u\in(\bigvee S)+\NN^d}x^{w\cdot u}
 =\frac{x^{w\cdot\bigvee S}}{\prod_j(1-x^{w_j})}.
\]
Apply finite inclusion--exclusion to Lemma~\ref{lem:dickson}.
Positive weights ensure that every coefficient counts only finitely
many vectors. Thus these are valid formal-power-series identities,
with no convergence issue. For $U=\varnothing$ the numerator is zero.
\end{proof}

\end{document}
